% !TeX program = xelatex
\documentclass[a4paper,11pt]{article}
\usepackage{kotex}
\usepackage[margin=23mm]{geometry}
\usepackage{booktabs,tabularx,array}
\usepackage{enumitem}
\usepackage{amsmath}
\usepackage{xcolor}
\usepackage[unicode,hidelinks]{hyperref}
\usepackage{xurl}
\usepackage{setspace}
\setstretch{1.15}
\setlength{\parindent}{0pt}
\setlength{\parskip}{5pt}
\setlength{\emergencystretch}{2em}
\renewcommand{\arraystretch}{1.2}
\newcolumntype{Y}{>{\raggedright\arraybackslash}X}
\setlist[enumerate]{leftmargin=1.6em,itemsep=3pt,topsep=3pt}
\setlist[itemize]{leftmargin=1.5em,itemsep=3pt,topsep=3pt}
\newcommand{\task}[3]{\subsection{#1}\textbf{작업 성격:} #2\par #3}
\hypersetup{pdftitle={보고서 및 데모 제출을 위한 추가 작업 계획}}
\begin{document}
\begin{center}
{\LARGE\bfseries 보고서 및 데모 제출을 위한 추가 작업 계획\par}
\vspace{6pt}
{\large Physical AI 작업 수행 검증 및 회귀 테스트 자동화\par}
\vspace{4pt}
{\small 기준일: 2026년 10월 20일 / 보고서 초안 v0.1 기준}
\end{center}

\section{목적과 현재 상태}
이 문서는 초기 기획서와 현재 구현·개발 실행 결과를 기준으로 남은 작업을 정리한다.
경량 AI 예측과 RGB-D 처리는 이번 작성 범위에서 제외한다. 우선 산출물은
A4 10쪽 보고서와 약 10분의 데모 영상이다. 아래 항목은 실행 계획이며, 작성만으로
추가 구현이나 실험이 완료된 것은 아니다.

현재 환경 후보 생성, 행동 피드백 탐색, 작업 목표 평가, 실패 증거 저장과 회귀
실행기가 구현되어 있다. 남은 작업의 상당 부분은 증거 정리와 실제 검증이다.
최신 AFS의 Random 비교 연결 등은 기존 코드의 재사용 가능성을 먼저 확인한 뒤
필요한 범위만 확장한다.

{\small
\begin{tabularx}{\linewidth}{@{}p{26mm}YY@{}}
\toprule
우선순위 & 작업 묶음 & 완료 후 확보할 결과 \\
\midrule
P0 / 제출 자료 & 성과 기준본·사례 근거·보고서·영상 & 수치와 주장, 실제 자료가 일치하는 제출본 \\
P1 / 핵심 검증 & 대표 사례의 회귀 테스트 & 저장 사례의 자동 재실행과 baseline 비교 증거 \\
P2 / 성능 평가 & 독립 평가·현재 AFS와 Random 비교 & 초기 정량 목표의 달성 여부를 판단할 자료 \\
후속 연구 & 추가 유형·휴대 가능한 자산·경계 검증 & 관측·재사용 범위 확대 \\
\bottomrule
\end{tabularx}}

P0는 현재 결과를 정직하게 보고하기 위한 필수 작업이다. P1은 과제명에 포함된
회귀 테스트의 실제 작동을 보여주기 위한 우선 검증이다. P2는 초기 성능 목표의
달성을 주장하려면 필요하다. 제출 평가에서 실제 목표 달성을 요구한다면 P2 역시
필수이며 문서 정리로 대체할 수 없다.

\section{P0: 제출 전 필수 자료 보완}

\task{확장 12종 결과와 초기 6종 비교표 확정}{기존 결과 재집계·문서화. 신규 로봇 실행 불필요.}{
초기 기획서의 다양성 목표는 충돌, 도달 불가, 장애물 간섭, 목표 위치 점유,
작업자 안전 위험, 인식 오류의 6종 중 4종 발견이다. 본 보고서는 실제 관찰에
따른 12종 확장을 주요 분석으로 채택하고, 초기 기준 결과를 비교용으로 병기한다.
12종 정의와 5종 관측을 구분한다. 최신 trajectory 스냅샷의 기존 3종 중 초기 기준에
포함되는 것은 목표 위치 점유와 작업자 안전 위험의 2종이다.}

\begin{itemize}
\item 실험별 시도·유효·PASS·FAIL·제외, 유형과 미분류 수를 원본 요약에서 확정한다.
\item 확장 12종의 지원·실제 관측·온라인 연결 여부와 초기 6종 결과를 각각 표시한다.
\item 추가 유형의 정의 계기, 검출 조건, 제외 조건과 실제 사례를 연결한다.
\item FDR의 분모는 유효 실행 수로 두고 전체 시도 대비 비율과 구분한다.
\item 다른 조건의 실험, 참고 이력, 신규 실행과 사후 재분류를 합산하지 않는다.
\end{itemize}
\textbf{완료 기준:} 보고서·도표·영상 대본이 하나의 기준 결과표를 사용하고,
모든 수치에 원본 파일과 분류 버전이 연결된다.

\task{실험 조건과 계획 변경 근거 정리}{원본 protocol·개발 기록 검토.}{
로봇 및 모델 식별자, 작업 목표와 완료 조건, 호출·시간 예산, 장면 스키마,
탐색 설정, 코드·자산 버전을 실험별로 정리한다. 초기 Panda/PyBullet 계획과
현재 주요 G1/MuJoCo 실험의 차이, 전환 시점과 이유는 기록을 근거로 기술한다.
근거가 없는 전환 이유를 사후에 만들어 쓰지 않는다.}

\textbf{완료 기준:} 보고서 2쪽의 계획 대비 표와 7쪽의 실험 조건표가 채워지고,
개발 중간 스냅샷과 완료된 실행을 구분한다.

\task{실제 행동 피드백 순환 1건 추적}{저장된 요청·응답·선택 기록의 연결.}{
이전 실행에서 무엇을 관찰했는지, LLM이 어떤 가설과 환경 변경을 제안했는지,
코드가 무엇을 선택했는지, 다음 실행에서 실제로 무엇이 측정되었는지를 연결한다.
예상이 틀린 경우도 그대로 보존한다.}

\textbf{완료 기준:} 이전 evidence ID, proposal ID, 기준 장면, 변경 축·값,
선택 근거, 다음 rollout ID와 결과가 실제 파일로 연결된 표 1개를 확보한다.
보고서 5쪽과 영상의 피드백 설명 구간에서 같은 자료를 사용한다.

\task{정상 예시와 대표 실패 사례 묶음}{영상 재생·로그 검토·캡처·편집.}{
아래 사례를 우선 후보로 검토한다. 최종 선택은 판정 근거의 완전성과 영상의
설명력을 확인한 뒤 확정한다. 서로 다른 조건의 정상·실패 사례를 엄밀한 대조
실험으로 설명하지 않는다.}

{\small
\begin{tabularx}{\linewidth}{@{}p{28mm}p{43mm}Y@{}}
\toprule
용도 & 후보 & 확인할 핵심 근거 \\
\midrule
정상 완료 & breadth / attempt\_00002 & 목표 도착·체류 등 완료 계약 \\
작업자 안전 위험 & trajectory / attempt\_00000 & 사람 proxy·노출 구간·분류 규칙 \\
물체 조작 실패 & trajectory / attempt\_00003 또는 00004 & 조작 요청·거절과 실제 실행 여부 \\
목표 위치 점유 & trajectory / attempt\_00005 & 종료 구간의 목표 영역 점유·정체 \\
\bottomrule
\end{tabularx}}

\textbf{완료 기준:} 사례마다 장면 이미지, 환경 변경값, 영상 타임코드,
행동·상태 근거, 최종 판정, 규칙 버전과 해석 한계를 묶는다. MP4의 존재·일부
메타데이터는 확인했지만 전체 재생과 정확한 편집 구간 선정은 남아 있다.

\task{보고서 PDF와 10분 데모 완성}{문서 편집·빌드·촬영·검수.}{
보고서의 보완 표시를 실제 자료로 채우고 XeLaTeX로 PDF를 빌드한다.
현재 작성 환경에는 LaTeX 엔진이 없어 실제 10쪽 여부와 시각 검수는 미완료다.
컴파일 후 쪽 수, 표·그림 넘침, 한글·경로 표시와 자료 참조를 확인한다.}

영상은 문제·구조 2분, 행동 피드백 1.5분, 실행 사례 2분, 판정·기록 1분,
회귀 테스트 2분, 결과·한계 1.5분을 기준으로 총 10분을 구성한다.
추론 대기는 편집하되 배속·생략과 저장 영상 재생 여부를 표시한다.

\textbf{완료 기준:} PDF와 영상의 수치가 일치하고, 영상 전체를 재생해
자막·음성·링크·민감정보 노출 여부를 확인한다. 실제 회귀 결과가 없다면
회귀 구간은 계획·기능 소개로 표시한다.

\section{P1: 대표 사례의 실제 회귀 검증}
\task{기존 회귀 실행기로 사례 묶음 검증}{기존 기능 검증 우선. 호환 문제에 한해 추가 구현.}{
저장 사례의 재실행과 baseline 비교 실행기는 이미 있다. PASS 1건과 목표 점유·
안전 위험 FAIL 각 1건을 우선 묶음으로 선정하고 읽기 전용 계획 검사부터 수행한다.
장면·목표·예산·자산이 현재 실행기에 맞는지 확인한다.}

실제 검증의 예산안은 사례당 2회, 총 최대 6회 시도다. 각 사례의 원본 호출 상한이
10이면 로봇 정책 호출 상한은 60회이며, 실제 상한은 계획 검사 결과로 다시 계산한다.
AFS 요청은 0회다. 이 값은 제안이며 실행 승인을 받거나 실제 실행한 예산이 아니다.

\begin{itemize}
\item 사례별 원래 장면과 작업 목표·예산을 유지하고 현재 코드·자산 조건을 기록한다.
\item 제외도 시도 예산을 소비하며 성공할 때까지 자동 반복하지 않는다.
\item baseline과 신규 PASS/FAIL, 유지·개선·퇴행·혼합·판정 불가 상태를 비교한다.
\item 새 회귀 결과를 과거 탐색의 FDR·Gain 표본에 추가하지 않는다.
\end{itemize}
\textbf{완료 기준:} 실제 실행 기록과 영상, baseline 비교표가 보고서 9쪽에 연결된다.
소수 반복은 기능 시연 근거이며 통계적인 개선 증명이 아니다.

\section{P2: 초기 정량 목표를 위한 추가 평가}
\task{현재 AFS와 Random 비교 연결}{기존 비교 실행기 재사용 검토 및 필요한 확장.}{
기존 AFS/Random 캠페인 실행기는 있지만 최신 행동 피드백 AFS와 동일한 방법을
평가하는지 확인해야 한다. 기존 명령을 실행한 사실만으로 최신 방법의 Gain이
검증되는 것은 아니다. 비교 연결은 구현 검증과 실제 성능 평가를 나누어 진행한다.}

\begin{enumerate}
\item 현재 제출 대상 AFS의 제안·선택·피드백 정책을 고정한다.
\item 같은 장면 도메인·로봇·목표·규칙·유효 실행 예산의 Random 대조를 연결한다.
\item seed, cold start, 제외 처리, 시도 상한과 비용 집계를 사전 정의한다.
\item 예산·조건 불일치, 재개·중단, 저장 응답의 중복 사용 등을 오프라인 검사한다.
\item 별도 예산을 확정한 뒤 실제 실행하고 seed별 결과를 모두 보고한다.
\end{enumerate}

\textbf{완료 기준:} 동일 유효 예산의 FDR, 초기 6종 다양성, 제외율과 사용량을
나란히 제시할 수 있다. 상대 향상으로 정의할 경우
\[
\mathrm{Gain}=\frac{\mathrm{FDR}_{\mathrm{AFS}}-\mathrm{FDR}_{\mathrm{Random}}}
{\mathrm{FDR}_{\mathrm{Random}}}
\]
를 사용하며, Random FDR이 0이면 비율을 계산하지 않고 실패 수와 절대 차이를
보고한다. 정의와 예외 처리는 실행 전에 고정한다.

\task{확장 12종 체계의 독립 평가 설계}{측정·설계 보완 및 실제 실행.}{
현재 독립 평가 실행기는 외부·다른 seed 이력을 섞지 않는 구조를 제공하지만
그 자체로 Random 비교를 완료하지는 않는다. 현재 평가 설계는 10종 목표이므로
12종 체계로 평가하려면 분류 프로필·보고·설계 manifest를 함께 버전 관리해야 한다.
목표 유형 수와 예산은 새 실행 전에 정하며, 기존 표본의 5종 관측을 독립 평가로
간주하지 않는다. 초기 6종 기준 결과도 비교용으로 남긴다.}

현재 주요 G1 경로에서 초기 6종 중 충돌·장애물 간섭·목표 점유·안전 위험은
운영 규칙이 있고, 도달 불가·인식 오류는 미지원이다. 현 지원 범위에서 4종을
관측하려면 지원하는 네 유형을 모두 실제로 발견해야 한다. 대규모 실행 전에
각 규칙의 관측 가능성과 장면 도메인을 감사한다. 목표 달성을 위해 기준이나
분모를 결과에 맞춰 바꾸지 않는다.

\textbf{완료 기준:} 고정된 조건·예산·12종 프로필로 평가하고 사전 정의한 확장 목표의
달성 여부를 보고한다. 초기 목표인 FDR 30\%, Random Gain 20\%, 6종 중 4종과의
비교도 구분해 남긴다. 외부 제출 기준이 초기 계획을 요구하면 해당 기준을 따른다.
실제 평가를 수행하지 않으면 미검증·미달성 상태를 유지한다.

\section{후속 연구와 이번 제출의 선행 조건이 아닌 작업}
{\small
\begin{tabularx}{\linewidth}{@{}p{36mm}YY@{}}
\toprule
항목 & 현재 상태 & 후속 작업 \\
\midrule
새 12종 온라인 탐색 & 사후 분류를 보고서의 주 분석으로 채택 & 자동 탐색까지 주장할 경우 새 조건에서 피드백·schema·선택·보고 연결 필요 \\
휴대 가능한 테스트 자산 & 원본 archive·외부 자산에 의존 & 경로·자산 명세, 해시와 복원 절차 정리 \\
최소 변화 경계 & 주요 G1 결과의 검증 근거 부족 & 고정 조건 대조·반복과 경계 측정 계약 수립 \\
피드백 기여도 & 제거 대조 결과 없음 & 같은 예산에서 피드백 유무 비교 \\
추가 실패 유형 & 도달 불가·인식 오류 미지원 & 주장 필요 시 별도 관측 계약·검출·검증 \\
\bottomrule
\end{tabularx}}

인식 오류 계측은 이번 범위에서 제외한 처리와 연결될 수 있으므로 현 제출의
기본 구현 과제로 넣지 않는다. 해당 유형은 미지원으로 명시한다. 새 유형 연결이나
추가 로봇 행동 기술 개발을 문서·영상 완성의 선행 조건으로 두지 않는다.

\section{실행 순서와 제출 판단}
\begin{enumerate}
\item \textbf{근거 확정:} 12종 확장 근거·관측표와 초기 6종 비교표, 조건표, 실제 피드백 순환을 정리한다.
\item \textbf{사례·회귀 준비:} 정상·실패 사례를 선정하고 캡처·타임코드와 회귀 계획 검사를 완료한다.
\item \textbf{보고서·대본:} 초안의 보완 표시를 채우고 PDF와 촬영 대본을 검수한다.
\item \textbf{실행 검증:} 별도 예산에 따라 회귀 실행, 필요 시 정량 평가를 수행한다.
\item \textbf{최종 반영:} 실제 결과로 표·대본을 갱신하고 PDF·영상을 제출 형태로 정리한다.
\end{enumerate}

회귀 결과와 정량 평가가 추가되지 않아도 구현·개발 관측 중심 보고서는 작성할 수
있다. 다만 이를 초기 정량 목표의 완전 달성으로 표현할 수는 없다. 현재 계획에는
신규 유료 API·GPU 실행이나 과거 동결 캠페인의 재개가 포함되어 실행된 바 없다.

\section{관련 파일}
보고서 저장소: \path{/workspace/g1_failure/src/AI_Pro_Report}.
초기 기획서는 \path{init_report.md}, 본문 초안은 \path{report.tex}이다.
보고서 10쪽의 근거 목록에서 원본 실험 경로를 확인할 수 있다.

원본 구현은 \path{/workspace/g1_failure/src/physical_ai_test_gen/scene2test} 아래에 있다.
회귀 실행은 \path{tools/run_behavior_regression.py}, 비교 실행은
\path{tools/run_afs_benchmark.py}, 독립 평가는
\path{tools/run_discovery_evaluation.py}를 우선 검토한다.
이 문서는 과거 동결 실행의 코드·자산 검사를 우회하거나 그 예산을 초기화하는
지침이 아니다. 조건이 달라졌다면 별도의 실행 계획과 결과로 보존한다.
\end{document}
